\section*{Chapitre \ref{ch:b1:periodicity} --- Périodicité}

\begin{solution}{exo:b1:periodicity:1}
\ce{Ca} $[\ce{Ar}]\,4\text{s}^2$~: \omterm{def:b1:periodicity:table}{période} 4, \omterm{def:b1:periodicity:table}{groupe} 2, \omterm{def:b1:periodicity:table}{bloc} s. \ce{Fe}
$[\ce{Ar}]\,3\text{d}^6 4\text{s}^2$~: \omterm{def:b1:periodicity:table}{période} 4, \omterm{def:b1:periodicity:table}{groupe} 8, \omterm{def:b1:periodicity:table}{bloc} d. \ce{Se}
$[\ce{Ar}]\,3\text{d}^{10} 4\text{s}^2 4\text{p}^4$~: \omterm{def:b1:periodicity:table}{période} 4, \omterm{def:b1:periodicity:table}{groupe} 16,
\omterm{def:b1:periodicity:table}{bloc} p. \ce{I} $[\ce{Kr}]\,4\text{d}^{10} 5\text{s}^2 5\text{p}^5$~:
\omterm{def:b1:periodicity:table}{période} 5, \omterm{def:b1:periodicity:table}{groupe} 17, \omterm{def:b1:periodicity:table}{bloc} p.
\end{solution}

\begin{solution}{exo:b1:periodicity:2}
\ce{Na} est plus gros que \ce{Cl}~: même \omterm{def:b1:periodicity:table}{période}, $Z^{*}$ plus faible.
\ce{Na} est plus gros que \ce{Na+}, qui a perdu sa couche 3s. \ce{Cl-}
(\qty{181}{pm}) est plus gros que \ce{F-} (\qty{133}{pm})~: une couche de
plus. \ce{Na+} (\qty{102}{pm}) est plus gros que \ce{Mg^{2+}}
(\qty{72}{pm})~: mêmes dix électrons, charge nucléaire 11 contre 12.
% ledger: rion:Cl-VI, rion:F-VI, rion:Na+VI, rion:Mg2+VI
\end{solution}

\begin{solution}{exo:b1:periodicity:3}
$\ce{Na} (5.14) < \ce{Al} (5.99) < \ce{Mg} (7.65) < \ce{S} (10.36) <
\ce{P} (10.49) < \ce{Ar} (15.76)$. L'augmentation générale le long de la
\omterm{def:b1:periodicity:table}{période} est rompue deux fois~: \ce{Al} est sous \ce{Mg} (l'électron
arraché est le premier électron 3p, plus haut en énergie que 3s) et
\ce{S} sous \ce{P} (l'électron arraché est le premier électron 3p
apparié).
\end{solution}

\begin{solution}{exo:b1:periodicity:4}
La \omterm{def:b1:periodicity:polarisability}{polarisabilité} mesure la facilité avec laquelle le nuage électronique
se déforme sous l'effet d'un champ électrique. \ce{I-} est plus
polarisable que \ce{F-}~: il est plus gros et ses électrons externes sont
plus loin du noyau. \ce{Cl-} (18 électrons retenus par $Z = 17$) est
bien plus polarisable que \ce{Na+} (10 électrons retenus par $Z = 11$),
comme les anions le sont plus que les cations.
\end{solution}

\begin{solution}{exo:b1:periodicity:5}
Rapports~: $18.83/5.99 = 3.1$, $28.45/18.83 = 1.5$, $119.99/28.45 = 4.2$.
Le saut se produit après le troisième électron~: trois \omterm{def:b1:quantum-numbers:configuration}{électrons de
valence}, \omterm{def:b1:periodicity:table}{groupe} 13, \omterm{def:b1:periodicity:table}{période} 3~: c'est l'aluminium. Il forme \ce{Al^{3+}},
$1\text{s}^2 2\text{s}^2
2\text{p}^6$.
% ledger: ie:Al, ie2:Al, ie3:Al, ie4:Al
\end{solution}

\begin{solution}{exo:b1:periodicity:6}
L'électron ajouté à \ce{F-} entre dans la petite couche 2p, où il est
fortement repoussé par les sept \omterm{def:b1:quantum-numbers:configuration}{électrons de valence} déjà présents~;
dans la couche 3p, plus étendue, du chlore, la répulsion est plus faible
et l'énergie libérée plus grande. Dans l'azote ($2\text{p}^3$,
\omorbs{u,u,u}), l'électron supplémentaire devrait s'apparier dans une
\omterm{def:b1:quantum-numbers:orbital}{sous-couche} à moitié remplie~; dans le magnésium ($3\text{s}^2$), il
devrait inaugurer la \omterm{def:b1:quantum-numbers:orbital}{sous-couche} 3p, plus haute. Dans aucun des deux cas
l'anion n'est plus stable que l'atome et un électron libre.
\end{solution}

\begin{solution}{exo:b1:periodicity:7}
$\chi_M(\ce{Na}) = (5.14 + 0.55)/2 = \qty{2.85}{eV}$~; $\chi_M(\ce{Cl}) =
(12.97 + 3.61)/2 = \qty{8.29}{eV}$. L'écart est grand~: le doublet d'une
liaison \ce{Na-Cl} appartient presque entièrement au chlore, et le
chlorure de sodium est ionique, \ce{Na+ Cl-}.
\end{solution}

\begin{solution}{exo:b1:periodicity:8}
\ce{H-F}~: $\Delta\chi = 1.78$, à la limite du domaine ionique (en
pratique une liaison covalente très polarisée), \ce{F^{\delta-}}.
\ce{C-H}~: 0.35, quasi apolaire. \ce{Na-Cl}~: 2.23, ionique, \ce{Cl^{-}}.
\ce{C-Cl}~: 0.61, covalente polarisée, \ce{C^{\delta+}-Cl^{\delta-}}.
\ce{O-H}~: 1.24, covalente polarisée, \ce{O^{\delta-}-H^{\delta+}}.
\end{solution}

\begin{solution}{exo:b1:periodicity:9}
En descendant le groupe, l'électron de valence est dans une couche de
$n$ plus grand, plus loin du noyau et écranté par davantage d'\omterm{def:b1:quantum-numbers:configuration}{électrons
de cœur}~; il est arraché plus facilement. Le rubidium, une couche plus
loin, a une \omterm{def:b1:periodicity:ionisation-energy}{énergie de première ionisation} inférieure à \qty{4.34}{eV}.
\end{solution}

\begin{solution}{exo:b1:periodicity:10}
Le zinc est $[\ce{Ar}]\,3\text{d}^{10} 4\text{s}^2$~: son électron
provient d'une \omterm{def:b1:quantum-numbers:orbital}{sous-couche} 4s pleine~; le gallium,
$\dots 4\text{s}^2 4\text{p}^1$, perd son unique électron 4p, plus haut en énergie et
écranté par 4s et 3d pleines. L'arsenic, $4\text{p}^3$, perd un électron
d'une \omterm{def:b1:quantum-numbers:orbital}{sous-couche} à moitié remplie~; le sélénium, $4\text{p}^4$, perd
l'électron apparié, déstabilisé par la répulsion de son partenaire~: son
énergie d'ionisation est un peu plus faible.
\end{solution}

\begin{solution}{exo:b1:periodicity:11}
$\Delta E = (3.98 - 2.20)^2 = 1.78^2 = \qty{3.17}{eV} = 3.17 \times 96.5 =
\qty{306}{kJ/mol}$. D'où $D(\ce{H-F}) = \tfrac12(436 + 158) + 306 = 297 +
306 \approx \qty{603}{kJ/mol}$. Ce fort excédent est l'attraction
supplémentaire entre les charges partielles $\ce{H^{\delta+}}$ et
$\ce{F^{\delta-}}$, due à la grande différence d'\omterm{def:b1:periodicity:electronegativity}{électronégativité}.
\end{solution}

\begin{solution}{exo:b1:periodicity:12}
L'\omterm{def:b1:periodicity:electronegativity}{échelle de Pauling} est définie à partir des énergies des liaisons
A--B~: un élément qui ne forme aucune liaison n'a pas de valeur de
Pauling. L'hélium, le néon et l'argon ne forment aucun composé stable.
Le krypton et surtout le xénon ont des couches de valence plus grandes et
plus polarisables, et des énergies d'ionisation plus faibles
(\qty{14.0}{eV} pour le krypton contre \qty{15.76}{eV} pour l'argon)~;
ils sont oxydés par le fluor et l'oxygène et forment des composés comme
\ce{XeF2}, de sorte que leurs énergies de liaison peuvent être mesurées.
% ledger: ie:Kr, ie:Ar
\end{solution}

\begin{solution}{pb:b1:periodicity:1}
\textbf{1.} Chaque électron est arraché à un ion plus positif et plus
petit que le précédent, qui retient plus fortement ses électrons.
\textbf{2.} $15.04/7.65 = 1.97$~; $80.14/15.04 = 5.33$~; $109.27/80.14 = 1.36$.
\textbf{3.} Le saut se produit après deux électrons~: X a deux \omterm{def:b1:quantum-numbers:configuration}{électrons
de valence}, \omterm{def:b1:periodicity:table}{groupe} 2~; dans la \omterm{def:b1:periodicity:table}{période} 3, c'est le magnésium.
\textbf{4.} \ce{Mg^{2+}}, $1\text{s}^2 2\text{s}^2 2\text{p}^6$, la
configuration du néon.
\textbf{5.} Un métal~: \omterm{def:b1:periodicity:table}{groupe} 2, à gauche du tableau, avec de faibles
\omterm{def:b1:periodicity:ionisation-energy}{énergies de première ionisation}.
\textbf{6.} Le troisième électron provient du cœur 2p, de $n$ plus petit,
bien plus proche du noyau et à peine écranté, et il est arraché à un ion
doublement chargé.
% ledger: ie:Mg, ie2:Mg, ie3:Mg, ie4:Mg
\textbf{7.} Tous ont dix électrons~: $1\text{s}^2 2\text{s}^2 2\text{p}^6$.
\textbf{8.} Les mêmes dix électrons sont retenus par des charges
nucléaires 8, 9, 11, 12, 13~: le nuage se contracte à mesure que $Z$
augmente.
\textbf{9.} \ce{O^{2-}}~: le plus gros, avec la plus faible charge
nucléaire pour ses dix électrons, et c'est un anion.
\textbf{10.} $140/54 = 2.6$.
\textbf{11.} $198.8/2 = \qty{99.4}{pm}$~; l'anion est $181/99.4 = 1.8$
fois plus grand que le \omterm{def:b1:periodicity:radius}{rayon covalent}.
% ledger: re:Cl2, rion:Cl-VI
\textbf{12.} $\ce{Mg^{2+}} < \ce{Na+} < \ce{Cl-}$~: les deux cations
sont isoélectroniques (charge 12 contre 11), et \ce{Cl-} a une troisième
couche.
\textbf{13.} $\chi_M$ = 10.41 (\ce{F}), 8.29 (\ce{Cl}), 7.59 (\ce{Br}),
7.53 (\ce{O}), 6.26 (\ce{C}), en eV.
\textbf{14.} Mulliken~: $\ce{F} > \ce{Cl} > \ce{Br} > \ce{O} > \ce{C}$~;
Pauling~: $\ce{F} > \ce{O} > \ce{Cl} > \ce{Br} > \ce{C}$. L'oxygène passe
de la quatrième à la deuxième place.
\textbf{15.} $3.98/10.41 = 0.382$, $3.16/8.29 = 0.381$, $2.96/7.59 =
0.390$~: presque constant, environ 0.38~; pour les halogènes, les deux
échelles sont proportionnelles.
\textbf{16.} L'\omterm{def:b1:periodicity:electron-affinity}{affinité électronique} de l'oxygène est faible
(\qty{1.44}{eV}) parce que l'électron ajouté doit s'apparier dans la
couche 2p, compacte~; la valeur de Mulliken de l'atome libre sous-estime
l'attraction que l'oxygène exerce dans ses liaisons, que l'\omterm{def:b1:periodicity:electronegativity}{échelle de
Pauling}, construite à partir des liaisons, mesure directement.
\textbf{17.} Le chlore dans \ce{HCl}~; le fluor dans \ce{ClF}.
\textbf{18.} $\Delta\chi = 0.96$~: covalente polarisée.
\textbf{19.} $D(\ce{H-H}) = 2 \times 218.0 = \qty{436.0}{kJ/mol}$.
\textbf{20.} $D(\ce{Cl-Cl}) = 2 \times 121.3 = \qty{242.6}{kJ/mol}$.
\textbf{21.} $D(\ce{H-Cl}) = 218.0 + 121.3 - (-92.3) = \qty{431.6}{kJ/mol}$.
\textbf{22.} $\Delta E = 431.6 - \tfrac12(436.0 + 242.6) = 431.6 - 339.3 =
\qty{92.3}{kJ/mol}$. En symboles, $D(\ce{H-Cl}) = \Delta_fH(\ce{H}) +
\Delta_fH(\ce{Cl}) - \Delta_fH(\ce{HCl})$ tandis que la moyenne des
liaisons symétriques vaut $\Delta_fH(\ce{H}) + \Delta_fH(\ce{Cl})$, donc
$\Delta E =
-\Delta_fH(\ce{HCl})$.
\textbf{23.} $92.3/96.485 = \qty{0.957}{eV}$.
\textbf{24.} $|\chi(\ce{Cl}) - \chi(\ce{H})| = \sqrt{0.957} =
\textbf{0.98}$, contre $3.16 - 2.20 = 0.96$ dans les tables~:
l'arithmétique de Pauling, appliquée à des données modernes, redonne
la différence tabulée à 0.02 près.
\end{solution}
